% =====================================================================
% 1. Introdução
% Responsável: Pessoa 1 (Maria Eduarda)
%
% O que cobrir (cerca de 1 página):
% - Problema: LLMs alucinam, falham em raciocínio multi-hop e não
%   mostram de onde tiraram a resposta.
% - Por que grafos ajudam: conhecimento explícito, travessia de
%   relações, procedência verificável.
% - O que este trabalho faz: sistema GraphRAG/KAG sobre a história da
%   computação, com Neo4j + LLM.
% - Contribuições, em lista: (i) dataset, (ii) modelo LPG, (iii)
%   pipeline Text-to-Graph, (iv) recuperação híbrida multi-hop,
%   (v) interface com procedência, (vi) avaliação comparativa.
% - Parágrafo final com a organização do artigo (seção 2 trata de...).
% =====================================================================

\section{Introdução}

\pendente{introdução -- escrever depois que as seções técnicas estiverem prontas}
